\newcommand{\DoNotLoadEpstopdf}{}
\documentclass[11pt,reqno]{amsart}
\usepackage[T1]{fontenc}
\usepackage{lmodern,amsmath,amssymb,amsthm,mathtools,mathrsfs,microtype}
\usepackage[margin=1.1in]{geometry}
\usepackage[hidelinks]{hyperref}
\hypersetup{pdftitle={Collision records, arithmetic, and raw local inversion in a triangular Lorentz gas},pdfauthor={Qian Qi}}
\setlength{\emergencystretch}{2em}
\allowdisplaybreaks[2]
\numberwithin{equation}{section}
\newtheorem{theorem}{Theorem}[section]
\newtheorem{lemma}[theorem]{Lemma}
\newtheorem{proposition}[theorem]{Proposition}
\newtheorem{corollary}[theorem]{Corollary}
\theoremstyle{definition}
\newtheorem{definition}[theorem]{Definition}
\newtheorem{inputstatement}[theorem]{Input statement}
\theoremstyle{remark}\newtheorem{remark}[theorem]{Remark}
\newcommand{\R}{\mathbb R}\newcommand{\Z}{\mathbb Z}
\newcommand{\one}{\mathbf1}\newcommand{\dd}{\,\mathrm d}
\newcommand{\esssup}{\operatorname*{ess\,sup}}
\newcommand{\diver}{\operatorname{div}}\newcommand{\TV}{\operatorname{TV}}
\title[Collision records and raw local inversion]{Collision records, arithmetic, and raw local inversion in a triangular Lorentz gas}
\author{Qian Qi}
\date{October 13, 2026. A2-DYN, revision 99}
\subjclass[2020]{37D50, 37A50, 60F05, 28A75}
\keywords{Lorentz gas, exact collision records, endpoint momentum uniqueness, stable Jacobian, marked pressure sums, optical capacity, paired physical traces}
\begin{document}
\raggedbottom
\begin{abstract}
We study the exact-label pointwise raw local limit problem for actual
returns in a triangular Lorentz gas. Convexity of the Cartesian flight
action proves global uniqueness for a prescribed lifted itinerary and
two endpoint momenta. A comparison of the inverse optical twist with
the invariant stable Jacobian then controls the actual weighted
momentum-fiber multiplicity by a marked pressure sum. The mark removes
one stable factor; a finite physical step alphabet pays the resulting
multiplicity without an infinite homogeneity-strip count. Established
fixed-table pressure bounds imply bounded optical capacity at every
fixed radius. Thus the first inverse-gradient insertion, long critical
residual height and grazing-ended short-source estimates no longer
require a capacity hypothesis at fixed radius. A submultiplicative
finite-block theorem separately gives a radius-uniform exponential
growth bound and a sufficient spectral margin. The original uniform
pointwise target is unchanged: uniform certification of that margin
and decay of the positive nongrazing paired short trace remain to be
proved. All earlier proofs and both historical source variants are
retained.
\end{abstract}
\maketitle
\input{journal-problem-v93}

\subsection{The weighted multiplicity and the pressure comparison}
The principal target is \eqref{eq:v76-target}, with the original source,
normalization and complete arithmetic reference. The finite-count
quadratic cap and both physical residual mechanisms are retained below
with their proofs. Their capacity condition concerns a weighted sum over
\emph{endpoint-momentum fibers}, not an ordinary symbolic word count.
Theorem~\ref{thm:v99-unique} removes the possible multiplicity inside
one physical itinerary by a global Cartesian convexity argument.
Lemma~\ref{lem:v99-stable-weight} then compares its precise optical
weight with the stable Jacobian and removes each singular incidence
at its marked time.

The resulting theorem compares the complete momentum-fiber capacity
with geometric partition sums. For each fixed radius it is bounded
uniformly in collision count by the established fixed-table pressure
theory; see Corollary~\ref{cor:v99-fixed-radius}. For the radius-uniform
problem, Theorem~\ref{thm:v99-block} gives a finite-block exponential
growth bound and its explicit strict spectral margin. These are
different quantifiers. The former is not silently promoted to the
latter by compactness of the parameter interval.

At fixed radius the new estimate makes the critical recovery and the
complete short-trace comparison unconditional relative to the retained
analytic and geometric inputs. The positive same-segment minimum on
short nongrazing segments remains an additional source-height question,
even at one radius. The original uniform theorem is not replaced by
the fixed-radius consequence, nor is the paired minimum identified
with a projected signed cancellation.

\input{journal-inputs-v94}
\input{journal-residual-v98}
\input{core/235_quadratic_gradient_cap_and_uncut_insertion}
\input{core/236_disjoint_residuals_and_critical_height_certificate}
\input{core/237_paired_physical_traces_and_scalar_reduction}

\input{core/238_global_momentum_uniqueness_and_stable_weight}
\input{core/239_marked_pressure_comparison_and_growth}
\input{core/240_fixed_radius_recovery_and_uniform_scope}

\appendix
\input{appendices/v99_dependency_map}
\input{journal-references-v99}
\end{document}
